\documentclass[11pt,a4paper]{article}
\usepackage[margin=1in]{geometry}
\usepackage{amsmath,amssymb,amsthm}
\usepackage{algorithm}
\usepackage{algpseudocode}
\usepackage{booktabs}
\usepackage{hyperref}

\theoremstyle{definition}
\newtheorem{definition}{Definition}[section]
\newtheorem{theorem}{Theorem}[section]
\newtheorem{lemma}[theorem]{Lemma}
\newtheorem{remark}{Remark}[section]

\title{\textbf{Rigorous Mathematical Specification, Probabilistic Expectation Gating, and Gradient Analysis of Smooth Activations (GELU, SiLU/Swish, Mish)}}
\author{Team Scaibu \\ \small Email: \texttt{jaydeep.9664920749@gmail.com} \\ \small Architecture and Core Engine Team}
\date{October 2026}

\begin{document}
\maketitle

\begin{abstract}
This document presents the complete mathematical formulation and algorithmic specification of non-monotonic smooth activations: Gaussian Error Linear Unit (GELU exact and tanh approximation), Sigmoid Linear Unit (SiLU / Swish), and Mish. Unlike piecewise linear functions, smooth non-monotonic activations possess non-zero negative curvature and continuous higher-order derivatives, stabilizing optimization in modern Transformer (BERT, GPT, Llama) and Vision architectures. We formalize the stochastic regularization derivation, formulate complete algorithmic pseudo-code, derive complexity bounds, calculate a numerical example, and discuss Taylor series stability.
\end{abstract}

\section{Notation and Mathematical Preliminaries}

Let $X \in \mathbb{R}^{N \times D}$ denote a 2D pre-activation tensor.

\begin{itemize}
    \item $x \in \mathbb{R}$: Scalar pre-activation coordinate.
    \item $\Phi(x) \triangleq \frac{1}{\sqrt{2\pi}} \int_{-\infty}^x e^{-t^2/2} dt = \frac{1}{2}\left[1 + \operatorname{erf}\left(\frac{x}{\sqrt{2}}\right)\right]$: Standard Gaussian cumulative distribution function.
    \item $\sigma(x) \triangleq \frac{1}{1 + e^{-x}}$: Standard logistic sigmoid function.
    \item $\operatorname{softplus}(x) \triangleq \ln(1 + e^x)$: Smooth positive rectifier.
    \item $\phi(x) \in \mathbb{R}$: Output activation value.
    \item $\phi'(x)$: Exact analytical derivative.
\end{itemize}

\begin{table}[h]
\centering
\caption{Summary of Mathematical Symbols for Smooth Activations}
\begin{tabular}{lll}
\toprule
\textbf{Symbol} & \textbf{Type / Domain} & \textbf{Description} \\
\midrule
$x$ & $\mathbb{R}$ & Pre-activation scalar \\
$\Phi(x)$ & $[0, 1]$ & Standard Normal CDF \\
$\sigma(x)$ & $(0, 1)$ & Logistic sigmoid \\
$\phi(x)$ & $\mathbb{R}$ & Post-activation output \\
$\phi'(x)$ & $\mathbb{R}$ & Analytic derivative \\
\bottomrule
\end{tabular}
\end{table}

\section{Problem Definition and Core Concepts}

\subsection{Intuition and Motivation}
In standard ReLU ($\max(0, x)$), input values are deterministically multiplied by 0 or 1 based on sign alone. GELU bridges deterministic activation with stochastic dropout: an input $x$ is multiplied by $m \sim \operatorname{Bernoulli}(\Phi(x))$, where the probability that $x$ drops to 0 increases smoothly as $x$ becomes more negative. The expected activation is $\mathbb{E}[m \cdot x] = x \cdot \Phi(x)$. This continuous curve allows small negative values to propagate small negative gradients, preventing dead neurons.

\subsection{Formal Mathematical Definition}
\begin{definition}[Smooth Activation Equations and Derivatives]
\begin{enumerate}
    \item \textbf{GELU Exact}:
    \begin{align}
    \operatorname{GELU}(x) &= x \Phi(x) = \frac{x}{2} \left[ 1 + \operatorname{erf}\left( \frac{x}{\sqrt{2}} \right) \right] \label{eq:gelu_exact} \\
    \frac{d}{dx}\operatorname{GELU}(x) &= \Phi(x) + x \cdot \frac{e^{-x^2/2}}{\sqrt{2\pi}} \label{eq:gelu_grad}
    \end{align}
    \item \textbf{GELU Tanh Approximation}:
    \begin{align}
    \operatorname{GELU}_{\text{tanh}}(x) &= \frac{x}{2} \left[ 1 + \tanh\left( \sqrt{\frac{2}{\pi}} \left( x + 0.044715 x^3 \right) \right) \right] \label{eq:gelu_tanh}
    \end{align}
    \item \textbf{SiLU / Swish}:
    \begin{align}
    \operatorname{SiLU}(x) &= x \cdot \sigma(x) = \frac{x}{1 + e^{-x}} \label{eq:silu_def} \\
    \frac{d}{dx}\operatorname{SiLU}(x) &= \sigma(x) + x \sigma(x)(1 - \sigma(x)) = \operatorname{SiLU}(x) + \sigma(x)(1 - \operatorname{SiLU}(x)) \label{eq:silu_grad}
    \end{align}
    \item \textbf{Mish}:
    \begin{align}
    \operatorname{Mish}(x) &= x \cdot \tanh(\operatorname{softplus}(x)) = x \cdot \tanh(\ln(1 + e^x)) \label{eq:mish_def}
    \end{align}
\end{enumerate}
\end{definition}

\section{Algorithmic Specification}

The computational pipeline implemented in \texttt{NnAlgoSmoothActivations} is shown below.

\begin{algorithm}[H]
\caption{Smooth Activations Forward Evaluation and Analytic Derivatives}
\begin{algorithmic}[1]
\Require Input tensor $X \in \mathbb{R}^{N \times D}$, variant $\in \{\text{gelu\_exact}, \text{gelu\_tanh}, \text{silu}, \text{mish}\}$
\Ensure Activated tensor $Y \in \mathbb{R}^{N \times D}$, derivative tensor $G \in \mathbb{R}^{N \times D}$
\State $N \gets \text{rows}(X), \quad D \gets \text{cols}(X)$
\State $Y \gets \text{new Matrix of shape } (N, D), \quad G \gets \text{new Matrix of shape } (N, D)$
\State $C \gets \sqrt{2.0 / \pi} \approx 0.7978845608028654$
\For{$i \gets 1 \textbf{ to } N$}
    \For{$j \gets 1 \textbf{ to } D$}
        \State $x \gets X[i][j]$
        \If{$\text{variant} = \text{gelu\_exact}$}
            \State $\text{cdf} \gets 0.5 \cdot (1.0 + \operatorname{erf}(x / \sqrt{2.0}))$
            \State $Y[i][j] \gets x \cdot \text{cdf}$
            \State $G[i][j] \gets \text{cdf} + x \cdot (\exp(-0.5 \cdot x^2) / \sqrt{2.0 \cdot \pi})$
        \ElsIf{$\text{variant} = \text{gelu\_tanh}$}
            \State $u \gets C \cdot (x + 0.044715 \cdot x^3)$
            \State $t \gets \tanh(u)$
            \State $Y[i][j] \gets 0.5 \cdot x \cdot (1.0 + t)$
            \State $G[i][j] \gets 0.5 \cdot (1.0 + t) + 0.5 \cdot x \cdot (1.0 - t^2) \cdot C \cdot (1.0 + 3.0 \cdot 0.044715 \cdot x^2)$
        \ElsIf{$\text{variant} = \text{silu}$}
            \State $s \gets 1.0 / (1.0 + \exp(-x))$
            \State $Y[i][j] \gets x \cdot s$
            \State $G[i][j] \gets s \cdot (1.0 + x \cdot (1.0 - s))$
        \ElsIf{$\text{variant} = \text{mish}$}
            \State $\text{sp} \gets \ln(1.0 + \exp(x))$
            \State $t \gets \tanh(\text{sp})$
            \State $Y[i][j] \gets x \cdot t$
            \State $G[i][j] \gets t + x \cdot (1.0 - t^2) \cdot (1.0 / (1.0 + \exp(-x)))$
        \EndIf
    \EndFor
\EndFor
\State \Return $\{\text{output\_tensor}: Y, \text{gradients}: G\}$
\end{algorithmic}
\end{algorithm}

\section{Theoretical Guarantees and Smooth Curvature}

\begin{theorem}[Non-Monotonic Global Minimum]
Both $\operatorname{GELU}(x)$ and $\operatorname{SiLU}(x)$ are non-monotonic $\mathcal{C}^\infty$ smooth functions with a single unique global minimum in the negative half-plane:
\begin{align}
\min_{x \in \mathbb{R}} \operatorname{GELU}(x) &\approx -0.1700 \quad \text{at } x \approx -0.7518 \\
\min_{x \in \mathbb{R}} \operatorname{SiLU}(x) &\approx -0.2785 \quad \text{at } x \approx -1.2785
\end{align}
\end{theorem}

\subsection{Limitations}
\begin{itemize}
    \item \textbf{Higher Transcendental Cost}: Evaluating $\operatorname{erf}$ or $\tanh$ requires $\approx 5\times$ to $10\times$ more cycles than a simple $\max(0, x)$ comparison.
\end{itemize}

\section{Complexity Analysis}

\subsection{Time and Space Complexity}
\begin{itemize}
    \item \textbf{Time}: $\Theta(ND)$ transcendental arithmetic operations.
    \item \textbf{Space}: $\Theta(ND)$ auxiliary memory for activation caching.
\end{itemize}

\section{Worked Numerical Example}

Consider $x = 1.0$ under SiLU and GELU exact.

\begin{enumerate}
    \item \textbf{SiLU ($x=1.0$)}:
    \begin{align}
    \sigma(1.0) &= \frac{1}{1 + e^{-1}} \approx 0.7310586 \\
    Y_{\text{SiLU}} &= 1.0 \times 0.7310586 = 0.7310586 \\
    G_{\text{SiLU}} &= 0.7310586 \times (1.0 + 1.0 \times (1.0 - 0.7310586)) \\
    &= 0.7310586 \times (1.0 + 0.2689414) = 0.7310586 \times 1.2689414 \approx 0.927670
    \end{align}
    \item \textbf{GELU Exact ($x=1.0$)}:
    \begin{align}
    \Phi(1.0) &= 0.5(1 + \operatorname{erf}(1/\sqrt{2})) \approx 0.8413447 \\
    Y_{\text{GELU}} &= 1.0 \times 0.8413447 = 0.8413447 \\
    G_{\text{GELU}} &= 0.8413447 + 1.0 \times \frac{e^{-0.5}}{\sqrt{2\pi}} = 0.8413447 + 0.2419707 \approx 1.083315
    \end{align}
\end{enumerate}

\section{Numerical Considerations and Edge Cases}

\begin{itemize}
    \item \textbf{Mish Softplus Overflow Protection}: For $x > 20$, $\operatorname{softplus}(x) \approx x$ and $\tanh(x) \approx 1$, preventing float exponent overflow.
\end{itemize}

\begin{thebibliography}{9}
\bibitem{hendrycks2016} D.~Hendrycks and K.~Gimpel, ``Gaussian error linear units (GELUs),'' \emph{arXiv preprint arXiv:1606.08415}, 2016.
\bibitem{ramachandran2017} P.~Ramachandran, B.~Zoph, and Q.~V.~Le, ``Searching for activation functions,'' \emph{arXiv preprint arXiv:1710.05941}, 2017.
\bibitem{misra2019} D.~Misra, ``Mish: A self regularized non-monotonic activation function,'' \emph{arXiv preprint arXiv:1908.08681}, 2019.
\end{thebibliography}

\end{document}
